\documentclass[11pt,a4paper]{report}
% Packages
\usepackage[utf8]{inputenc}
\usepackage[english]{babel}
\usepackage{geometry}
\geometry{a4paper, margin=2.5cm}
\usepackage{graphicx}
\usepackage{booktabs}
\usepackage{longtable}
\usepackage{array}
\usepackage{float}
\usepackage{caption}
\usepackage{enumitem}
\usepackage{xcolor}
\usepackage{tikz}
\usetikzlibrary{shapes.geometric, arrows, positioning}
\usepackage{hyperref}
% Configure Hyperref
\hypersetup{
    colorlinks=true,
    linkcolor=blue,
    filecolor=magenta,      
    urlcolor=cyan,
    pdftitle={IoT-Based Multi-Sensor Anti-Theft System for Kitchen Area},
    pdfpagemode=FullScreen,
}
% Define Colors
\definecolor{fptorange}{RGB}{242, 113, 37}
\definecolor{darkblue}{RGB}{0, 51, 102}
% TikZ style configurations
\tikzstyle{startstop} = [draw, rectangle, rounded corners, minimum width=3cm, minimum height=1cm, text centered, draw=black, fill=red!30]
\tikzstyle{io} = [draw, trapezium, trapezium left angle=70, trapezium right angle=110, minimum width=3cm, minimum height=1cm, text centered, draw=black, fill=blue!30]
\tikzstyle{process} = [draw, rectangle, minimum width=3cm, minimum height=1cm, text centered, text width=3cm, draw=black, fill=orange!30]
\tikzstyle{decision} = [draw, diamond, minimum width=3cm, minimum height=1cm, text centered, text width=2cm, draw=black, fill=green!30]
\tikzstyle{arrow} = [thick,->,>=stealth]
\begin{document}
% ==========================================
% TITLE PAGE
% ==========================================
\begin{titlepage}
    \begin{center}
        % FPT University Logo
        \begin{figure}[H]
            \centering
            % Safely check and load FPT University Logo from workspace
            \includegraphics[width=0.45\textwidth]{images.png}
        \end{figure}
        
        \vspace{1.5cm}
        
        \textcolor{darkblue}{\Huge\bfseries IoT-BASED MULTI-SENSOR ANTI-THEFT SYSTEM FOR KITCHEN AREA} \\
        
        \vspace{1.5cm}
        
        \textcolor{fptorange}{\Large\bfseries ACADEMIC PROJECT REPORT} \\
        \large Course: IOT102 -- Internet of Things \\
        \large Class: SE2034 \\
        
        \vspace{2cm}
        
        \textbf{Prepared by Group 6:} \\
        \vspace{0.3cm}
        \begin{table}[h]
            \centering
            \begin{tabular}{llll}
                \toprule
                \textbf{No.} & \textbf{Full Name} & \textbf{Student ID} & \textbf{Project Role} \\
                \midrule
                1 & Ngo Gia Long & SE190732 & Group Leader \\
                2 & Nguyen Nhat Anh & SE192077 & Hardware Engineer \\
                3 & Vo Tran Cong Danh & SE190824 & Software Developer \\
                4 & Nguyen An Vuong & SE190996 & Testing and Support \\
                \bottomrule
            \end{tabular}
        \end{table}
        
        \vfill
        
        \large FPT University, Ho Chi Minh Campus \\
        \large June 2026
    \end{center}
\end{titlepage}
% ==========================================
% ABSTRACT & FRONT MATTER
% ==========================================
\pagenumbering{roman}
\begin{abstract}
Residential security has evolved rapidly with the rise of the Internet of Things (IoT). Kitchen areas represent a critical point of vulnerability in residential properties, as they often contain external entrance points (windows and back doors) and high-value appliances. Traditional security systems rely on single-sensor solutions that suffer from high false-alarm rates or are easily bypassed. This report presents a prototype-level design and implementation of an IoT-Based Multi-Sensor Anti-Theft System for Kitchen Areas utilizing an ESP32-CAM AI Thinker OV2640 with an ESP32-CAM-MB programmer. By incorporating sensor fusion logic involving a Passive Infrared (PIR) sensor, an ultrasonic distance sensor, and a Light Dependent Resistor (LDR), the system achieves resilient intrusion detection while significantly reducing false alarms caused by household pets. Furthermore, a DS1307 Real-Time Clock (RTC) ensures automated security state recovery after power failure events. The system interfaces with the Arduino IoT Cloud platform to provide a role-based user experience using three separate dashboards tailored for children, parents, and administrative monitoring. Real-time emergency escalation is supported via Google Apps Script for automated email warnings and a Telegram Bot for instant messaging notifications. Practical testing confirms that the multi-sensor approach improves security robustness and tampering resistance, demonstrating a highly effective local-to-cloud security integration.
\end{abstract}
\tableofcontents
\listoffigures
\listoftables
\clearpage
\pagenumbering{arabic}
% ==========================================
% SECTION 1: INTRODUCTION
% ==========================================
\chapter{Introduction}
\section{Project Overview}
The modern residential environment is increasingly threatened by unauthorized intrusions, requiring robust and intelligent home security solutions. While broad property monitoring is common, specific areas within a home present unique security challenges. The kitchen, which typically features external access windows, back doors, and valuable electrical appliances, is a primary entry candidate for intruders. 
This project details the development of an Internet of Things (IoT) based smart home security system specifically designed to monitor, protect, and report security breaches in a kitchen environment. Controlled by the ESP32-CAM AI Thinker OV2640 microcontroller, the system integrates a variety of peripheral sensors (PIR, ultrasonic, and light intensity sensors) to detect intrusions. When an alarm condition is met, the system triggers local warnings (LED and buzzer) and leverages cloud and API services (Arduino IoT Cloud, Google Apps Script, and Telegram Bot) to update variables, display statuses, and dispatch email and instant messaging alerts.
\section{Motivation and Problem Statement}
Homeowners face significant challenges with conventional security products:
\begin{itemize}
    \item \textbf{High False-Alarm Rates (False Positives):} Simple motion detectors cannot differentiate between a human intruder and a small pet (e.g., a cat or dog) moving across the floor. Repeated false alarms cause users to ignore notifications or disable the system entirely.
    \item \textbf{Vulnerability to Power Loss:} Standard microcontrollers lose their runtime state during power outages. If an intruder cuts the power to the kitchen before entering, standard software security modes are cleared, leaving the system disarmed when power returns.
    \item \textbf{Tampering and Sabotage:} Intruders can easily disable a sensor by covering its lens or tape-blocking the light path. Single-sensor configurations are unable to detect when they have been blinded.
    \item \textbf{Monolithic Interfaces:} A single control dashboard is poorly suited for multi-generational families. Children require simple panic mechanisms without configuration controls, while parents require comprehensive override and arming options.
\end{itemize}
This project is motivated by the need to resolve these limitations by combining multi-sensor fusion, low-cost microcontrollers, and modern cloud notification channels.
\section{System Objectives}
The core objectives of the proposed system are:
\begin{enumerate}
    \item \textbf{Improve Detection Reliability:} Implement sensor fusion logic where PIR, ultrasonic, and light sensors validate each other to confirm actual human presence.
    \item \textbf{Reduce False Alarms:} Differentiate between humans and domestic animals using spatial verification.
    \item \textbf{Ensure System Resiliency:} Maintain security states across power losses utilizing an independent real-time clock.
    \item \textbf{Provide Tamper Prevention:} Identify physical attempts to blind or block sensors and trigger immediate high-priority alerts.
    \item \textbf{Deliver Role-Based UI:} Partition control access into child-friendly, parent-level, and admin-level views.
    \item \textbf{Enable Multi-Channel Notifications:} Deliver instant alerts via Arduino Cloud variables, Telegram messages, and structured emails.
\end{enumerate}
% ==========================================
% SECTION 2: SYSTEM ARCHITECTURE
% ==========================================
\chapter{System Architecture}
\section{Proposed System Architecture}
The system architecture follows a hierarchical structure consisting of physical sensors, a local processing core (ESP32-CAM), an IoT cloud middleware (Arduino Cloud), and third-party alert dispatchers (Google Apps Script and Telegram API). Figure \ref{fig:block_diagram} shows the high-level flow of the system.
\begin{figure}[H]
\centering
\begin{tikzpicture}[nodeDistance=2cm, auto]
    % Nodes
    \node [process, fill=blue!10] (inputs) {\textbf{Inputs} \\ PIR Sensor \\ Ultrasonic Sensor \\ LDR Light Sensor \\ DS1307 RTC \\ RFID RC522 \\ Cloud SOS Buttons};
    \node [process, right=1.5cm of inputs, minimum height=2cm, fill=orange!15] (esp) {\textbf{Microcontroller} \\ ESP32-CAM AI Thinker \\ (Logic Processing \& \\ WiFi Connectivity)};
    \node [process, right=1.5cm of esp, fill=green!10] (local_out) {\textbf{Local Outputs} \\ Red LED \\ Passive Buzzer};
    \node [process, below=1.5cm of esp, fill=yellow!15] (cloud) {\textbf{Arduino IoT Cloud} \\ Syncs Cloud Variables \\ Update Dashboards \\ (Child / Parent / Admin)};
    \node [process, below=1.5cm of cloud, fill=purple!10] (notif) {\textbf{Notification Services} \\ Google Apps Script \\ Telegram Bot};
    \node [process, right=1.5cm of notif, fill=red!10] (out_notif) {\textbf{Remote Alerts} \\ Gmail Notifications \\ Telegram Message};
    % Paths
    \draw [arrow] (inputs) -- (esp);
    \draw [arrow] (esp) -- (local_out);
    \draw [arrow] (esp) -- (cloud);
    \draw [arrow] (cloud) -- (notif);
    \draw [arrow] (notif) -- (out_notif);
    \draw [arrow] (cloud) -- (esp);
\end{tikzpicture}
\caption{System Block Diagram and Information Flow}
\label{fig:block_diagram}
\end{figure}
\section{Hardware Component Selection}
The chosen hardware provides a balanced trade-off between cost, power consumption, and processing capacity. Table \ref{tab:hardware} details the hardware specification used.
\begin{table}[H]
\centering
\caption{Hardware Components and Specifications}
\label{tab:hardware}
\begin{tabular}{lp{6cm}p{6cm}}
    \toprule
    \textbf{Component} & \textbf{Model Used} & \textbf{Role in the System} \\
    \midrule
    Microcontroller & ESP32-CAM AI Thinker OV2640 with ESP32-CAM-MB & Core microcontroller with integrated WiFi, Bluetooth, and camera. Processes sensor logic and syncs with cloud. \\
    Real-Time Clock & DS1307 RTC & Keeps date/time tracking active via external CR2032 battery backup during power losses. \\
    RFID Module & RFID RC522 & Facilitates physical user authentication for arming/disarming. \\
    Ultrasonic Sensor & HC-SR04 / HY-SRF05 & Measures distance of obstacles to confirm entry and filter pets. \\
    Motion Sensor & PIR HC-SR501 / PIR 5050 & Detects body heat motion in the monitored area. \\
    Light Sensor & Light Dependent Resistor (LDR) Module & Measures light intensity to identify flashlights or sensor blockage. \\
    Acoustic Alarm & Passive Buzzer 1407 & Generates high-frequency alarm tones locally. \\
    Visual Indicator & Red LED (with 220$\Omega$ resistor) & Provides visual status and warning cues. \\
    Accessories & Breadboard, jumper wires, resistors, power adapter & Interconnects modules and supplies 5V DC power. \\
    \bottomrule
\end{tabular}
\end{table}
\section{Hardware Interfacing and Pin Connections}
To handle the limited GPIO pins on the ESP32-CAM module (since many pins are reserved for the camera array and flash LED), the connections are designed to prevent conflicts. Table \ref{tab:connections} documents the pin mapping.
\begin{table}[H]
\centering
\caption{ESP32-CAM Pin Interface Matrix}
\label{tab:connections}
\begin{tabular}{lllp{6cm}}
    \toprule
    \textbf{Peripheral} & \textbf{Module Pin} & \textbf{ESP32-CAM Pin} & \textbf{Interface Type \& Notes} \\
    \midrule
    DS1307 RTC & SDA & GPIO14 & I2C Data \\
               & SCL & GPIO15 & I2C Clock \\
    RFID RC522 & SDA (SS) & GPIO13 & SPI Slave Select \\
               & SCK & GPIO12 & SPI Clock \\
               & MOSI & GPIO2 & SPI Master Out Slave In (Shared with Buzzer) \\
               & MISO & GPIO16 & SPI Master In Slave Out \\
               & RST & GPIO0 & Reset (Required for boot configuration) \\
    PIR Sensor & OUT & GPIO3 & Digital Input (Requires pull-down) \\
    Ultrasonic Sensor & TRIG & GPIO1 & Digital Output (Trigger pulse) \\
                      & ECHO & GPIO13 & Digital Input (Needs $1k\Omega/2k\Omega$ voltage divider) \\
    LDR Module & AO & GPIO33 & Analog Input (ADC pin) \\
    Buzzer 1407 & + & GPIO2 & Digital Output (High triggers tone) \\
    Red LED & Anode (+) & GPIO12 & Digital Output (Via $220\Omega$ resistor) \\
    \bottomrule
\end{tabular}
\end{table}
The ultrasonic Echo line outputs a 5V signal, which must be attenuated to a safe 3.3V logic level using a voltage divider before reaching ESP32-CAM GPIO13. The voltage divider is configured using a $1k\Omega$ resistor in series with the Echo line and a $2k\Omega$ resistor connected to ground.
% ==========================================
% SECTION 3: CLOUD AND SOFTWARE ARCHITECTURE
% ==========================================
\chapter{Cloud and Software Architecture}
\section{Single Thing and Multi-Dashboard Paradigm}
The software model matches a single physical device (ESP32-CAM board) with a single logical entity (a ``Thing'' named \texttt{IoT\_Anti\_Theft\_System}) in the Arduino IoT Cloud. This centralized approach guarantees variable consistency. Instead of splitting the logic into multiple micro-devices, the access levels are handled by separating the user interface into three distinct dashboards according to the role of the user. Table \ref{tab:dashboards} presents the dashboard segmentation.
\begin{table}[H]
\centering
\caption{Role-Based Dashboard Design}
\label{tab:dashboards}
\begin{tabular}{lp{4cm}p{7cm}}
    \toprule
    \textbf{Dashboard Name} & \textbf{Target Audience} & \textbf{Widget Components and Access Rights} \\
    \midrule
    Child SOS Dashboard & Children, elderly, or non-technical residents & Minimal layout. Features a large, high-visibility \texttt{sos\_child} trigger button and status fields (\texttt{alarm\_status}, \texttt{last\_event}). No system arming or reset toggles. \\
    Parent Control Dashboard & Primary property owners & Intermediate layout. Contains arming switches (\texttt{alarm\_enabled}, \texttt{system\_armed}), reset button (\texttt{reset\_alarm}), adult SOS switch (\texttt{sos\_adult}), email status indicator (\texttt{email\_sent\_status}), and sensor values. \\
    Admin Demo Dashboard & Systems administrators and demonstration evaluators & Complete layout. Provides full debugging controls, including \texttt{demo\_mode} and \texttt{demo\_scenario} selectors, plus direct status monitoring for every sensor value and output. \\
    \bottomrule
\end{tabular}
\end{table}
\section{Arduino Cloud Variables}
Data synchronization is handled through Cloud Variables, which are automatically mirrored between the ESP32-CAM software and the Arduino IoT Cloud. The complete list of cloud variables is documented in Table \ref{tab:variables}.
\begin{longtable}{lllp{2.5cm}p{4cm}}
\caption{Arduino IoT Cloud Variables Matrix} \label{tab:variables} \\
\toprule
\textbf{Variable Name} & \textbf{Data Type} & \textbf{Permission} & \textbf{Update Policy} & \textbf{Purpose} \\
\midrule
\endfirsthead
\multicolumn{5}{c}%
{{\bfseries Table \thetable\ table continued from previous page}} \\
\toprule
\textbf{Variable Name} & \textbf{Data Type} & \textbf{Permission} & \textbf{Update Policy} & \textbf{Purpose} \\
\midrule
\endhead
\bottomrule
\endfoot
\bottomrule
\endlastfoot
\texttt{alarm\_enabled} & \texttt{bool} & Read \& Write & On Change & User switch to request security activation. \\
\texttt{system\_armed} & \texttt{bool} & Read Only & On Change & Confirms if the system is currently actively armed. \\
\texttt{reset\_alarm} & \texttt{bool} & Read \& Write & On Change & Command to silence alarms and clear states. \\
\texttt{alarm\_status} & \texttt{String} & Read Only & On Change & Descriptive status (e.g., SAFE, ARMED, INTRUSION). \\
\texttt{sos\_child} & \texttt{bool} & Read \& Write & On Change & Trigger indicating child emergency request. \\
\texttt{sos\_adult} & \texttt{bool} & Read \& Write & On Change & Trigger indicating adult emergency request. \\
\texttt{sos\_level} & \texttt{int} & Read Only & On Change & Priority rating (1 to 3) computed by the board. \\
\texttt{sos\_message} & \texttt{String} & Read Only & On Change & Formatted text describing the active SOS situation. \\
\texttt{pir\_detected} & \texttt{bool} & Read Only & On Change & Reflects the current state of the PIR sensor. \\
\texttt{ultrasonic\_distance} & \texttt{float} & Read Only & Periodic (5s) & Current distance measurement from the ultrasonic sensor. \\
\texttt{object\_near} & \texttt{bool} & Read Only & On Change & Set to true when distance is less than 100 cm. \\
\texttt{pet\_detected} & \texttt{bool} & Read Only & On Change & Indicates an ignored small/pet-like target. \\
\texttt{ldr\_value} & \texttt{int} & Read Only & Periodic (5s) & Raw light level reading. \\
\texttt{ldr\_delta} & \texttt{int} & Read Only & Periodic (5s) & Calculated light level delta over a sliding time window. \\
\texttt{light\_abnormal} & \texttt{bool} & Read Only & On Change & Active when a sudden flashlight pattern is matched. \\
\texttt{ldr\_covered} & \texttt{bool} & Read Only & On Change & True when the LDR is blocked/blinded. \\
\texttt{intrusion\_alert} & \texttt{bool} & Read Only & On Change & Confirms that the intrusion criteria are met. \\
\texttt{intrusion\_score} & \texttt{int} & Read Only & On Change & Danger points calculated for decision logic. \\
\texttt{threat\_level} & \texttt{int} & Read Only & On Change & Scaled indicator from 0 (Safe) to 4 (Emergency). \\
\texttt{sabotage\_alert} & \texttt{bool} & Read Only & On Change & Active when a sensor cover attempt is detected. \\
\texttt{device\_tampered} & \texttt{bool} & Read Only & On Change & Tracks physical interference with the hardware. \\
\texttt{buzzer\_on} & \texttt{bool} & Read Only & On Change & Tracks the state of the physical buzzer output. \\
\texttt{led\_red\_on} & \texttt{bool} & Read Only & On Change & Tracks the state of the red LED output. \\
\texttt{send\_email\_request} & \texttt{bool} & Read Only & On Change & Triggers the external Google Apps Script. \\
\texttt{email\_event\_type} & \texttt{String} & Read Only & On Change & Identifies the warning type (e.g., intrusion, SOS). \\
\texttt{email\_sent\_status} & \texttt{String} & Read Only & On Change & Tracks status (IDLE, SENDING, SENT, FAILED). \\
\texttt{last\_event} & \texttt{String} & Read Only & On Change & Text log showing the latest registered event. \\
\texttt{last\_event\_type} & \texttt{String} & Read Only & On Change & Categorization code for the last event. \\
\texttt{event\_counter} & \texttt{int} & Read Only & On Change & Running count of logged warning events. \\
\texttt{cooldown\_active} & \texttt{bool} & Read Only & On Change & Prevents rapid, repetitive email alerts. \\
\texttt{last\_alert\_time} & \texttt{String} & Read Only & On Change & Timestamp of the last email alert sent. \\
\texttt{demo\_mode} & \texttt{bool} & Read \& Write & On Change & Allows scenario testing via software. \\
\texttt{demo\_scenario} & \texttt{int} & Read \& Write & On Change & Code selector for testing specific mock events. \\
\end{longtable}
\section{Development Tools and Cloud Services}
The integration of development environments, IoT backends, and notification scripts forms the functional backbone of the system.
\begin{itemize}
    \item \textbf{Arduino Cloud:} Functions as the primary central database. Rather than requiring the ESP32-CAM to run local web servers or parse complex configuration files, the board maintains a persistent WebSocket connection to Arduino Cloud. Changes made to user controls on the phone are pushed instantly to the ESP32-CAM, and sensor values are streamed back to the cloud database.
    \item \textbf{Arduino IoT Remote Mobile App:} The mobile dashboard client. Available for Android and iOS, it renders the widgets (buttons, gauges, charts, labels) configured on the web. By logging in with different credentials or configuring customized views, residents access either the Child, Parent, or Admin interfaces.
    \item \textbf{Google Apps Script:} Serves as the email automation engine. The ESP32-CAM lacks the processing overhead to easily execute TLS connections for SMTP servers. To bypass this, the board issues an HTTP POST request to a deployed Google Apps Script Web App URL containing JSON parameters. The script, running inside the Google Cloud sandbox, utilizes the \texttt{MailApp} library to format and dispatch warning emails to configured family contacts using Gmail.
    \item \textbf{Telegram Bot:} Provides real-time notifications. The Telegram Bot API allows the system to send warnings and alert messages to a dedicated chat group. By sending simple HTTPS requests, the ESP32-CAM delivers alerts directly to users' phones within seconds.
\end{itemize}
\subsection{Tool Integration Flow}
The interaction between components follows a sequential path:
\begin{center}
\fbox{Sensors or Arduino IoT Remote Input} $\rightarrow$ \fbox{ESP32-CAM Processing} $\rightarrow$ \fbox{Local LED/Buzzer Alarm} $\rightarrow$ \fbox{Arduino Cloud Variable Update} $\rightarrow$ \fbox{Google Apps Script Email Alert} $\rightarrow$ \fbox{Telegram Bot Instant Message}
\end{center}
This structure ensures that local alarms respond immediately, while remote notifications compile and send in parallel without blocking local processing loops.
% ==========================================
% SECTION 4: SYSTEM OPERATIONS AND LOGIC
% ==========================================
\chapter{System Operations and Logic}
\section{General Operation Flow}
Upon cold boot, the system executes an initialization sequence. It configures local GPIO pins, starts I2C communication with the DS1307 RTC, reads the initial time, attempts a connection to the local WiFi router, and establishes a connection with the Arduino IoT Cloud Thing. Once online, the program enters a continuous monitoring loop, checking sensor values, comparing current time against the armed schedule, and evaluating cloud commands. Figure \ref{fig:workflow} illustrates the general system flow.
\begin{figure}[H]
\centering
\begin{tikzpicture}[nodeDistance=2.0cm, auto]
    % Flowchart Nodes
    \node [startstop] (start) {Start / Boot};
    \node [process, below=1.2cm of start] (init) {Initialize GPIO, RTC, WiFi \& Cloud};
    \node [process, below=1.2cm of init] (read) {Read Sensors (PIR, LDR, Ultrasonic, DS1307)};
    \node [decision, below=1.5cm of read] (armed) {System Armed?};
    \node [process, right=2.0cm of armed] (idle) {Log Low-Risk Events \& Maintain Sync};
    \node [process, below=1.5cm of armed] (eval) {Evaluate Fusion \& Tampering Logic};
    \node [decision, below=1.5cm of eval] (alert) {Alarm Condition?};
    \node [process, below=1.5cm of alert] (trigger) {Activate LED/Buzzer, Push Cloud Variables, Dispatch Email \& Telegram};
    \node [process, left=2.0cm of trigger] (cooldown) {Enter Cooldown \& Await Reset Command};
    
    % Connections
    \draw [arrow] (start) -- (init);
    \draw [arrow] (init) -- (read);
    \draw [arrow] (read) -- (armed);
    \draw [arrow] (armed) -- node [near start] {No} (idle);
    \draw [arrow] (armed) -- node {Yes} (eval);
    \draw [arrow] (idle) |- (read);
    \draw [arrow] (eval) -- (read); % Default feedback loop
    \draw [arrow] (eval) -- (alert);
    \draw [arrow] (alert) -- node [near start] {Yes} (trigger);
    \draw [arrow] (alert) -| node [near start] {No} (read);
    \draw [arrow] (trigger) -- (cooldown);
    \draw [arrow] (cooldown) |- (read);
\end{tikzpicture}
\caption{System Operation Flowchart}
\label{fig:workflow}
\end{figure}
\section{Detailed Function Explanations}
\subsection{Multi-Sensor Intrusion Detection}
\begin{itemize}
    \item \textbf{Real-Life Scenario:} An intruder opens a kitchen window at 2:00 AM, using a flashlight to scan the room.
    \item \textbf{Sensors Involved:} PIR motion sensor, HC-SR04/HY-SRF05 ultrasonic sensor, and LDR light sensor.
    \item \textbf{Data Read:} PIR digital output (\texttt{pir\_detected}), ultrasonic distance in cm (\texttt{ultrasonic\_distance}), and raw LDR light intensity (\texttt{ldr\_value}).
    \item \textbf{Decision Logic:} When the system is armed, it checks motion, distance, and light conditions. If the PIR detects motion, the ultrasonic sensor confirms an object is close to the protected area (distance is under 100 cm), and the LDR reports an abnormal light change (light level increases rapidly due to a flashlight beam), the system confirms a threat.
    \item \textbf{Local Alarm:} Activates the red LED and sounds the passive buzzer.
    \item \textbf{Cloud/Email/Telegram Notification:} Updates the cloud status to \texttt{INTRUSION\_ALERT}, logs the event time, sends an email alert via Google Apps Script, and dispatches a Telegram warning message.
    \item \textbf{Dashboard View:} The parent and admin dashboards display a red alert status, a threat level of 3, and log the event as \texttt{INTRUSION\_ALERT}.
    \item \textbf{Limitation/Assumption:} Assumes the kitchen is dark by default, making flashlight beams easily distinguishable.
\end{itemize}
\subsection{Pet False-Alarm Filtering}
\begin{itemize}
    \item \textbf{Real-Life Scenario:} A household cat jumps onto the kitchen counter or walks across the kitchen floor at night.
    \item \textbf{Sensors Involved:} PIR motion sensor and HC-SR04/HY-SRF05 ultrasonic sensor.
    \item \textbf{Data Read:} PIR state (\texttt{pir\_detected}) and distance (\texttt{ultrasonic\_distance}).
    \item \textbf{Decision Logic:} The PIR sensor has a wide field of view and will detect motion from both humans and pets. However, the ultrasonic sensor, mounted at a height of 80 cm parallel to the floor, measures the height and depth profile of moving objects. If motion is detected but the object is measured at a low height or is outside the narrow detection zone, the system classifies it as a pet.
    \item \textbf{Local Alarm:} Remains silent. No LED or buzzer is triggered.
    \item \textbf{Cloud/Email/Telegram Notification:} No high-priority alerts are sent. The event is logged in the cloud database as a low-risk log (\texttt{PET\_IGNORED}).
    \item \textbf{Dashboard View:} Displays status \texttt{SAFE} or \texttt{PET\_IGNORED} on the parent and admin dashboards.
    \item \textbf{Limitation/Assumption:} Assumes the pet is significantly smaller and lower than a human intruder. Large dogs may occasionally trigger the alarm if they jump high.
\end{itemize}
\subsection{Power-Loss Recovery using DS1307 RTC}
\begin{itemize}
    \item \textbf{Real-Life Scenario:} An intruder cuts the main power lines to the kitchen area to disable the security system.
    \item \textbf{Sensors Involved:} DS1307 Real-Time Clock with CR2032 backup battery, ESP32-CAM.
    \item \textbf{Data Read:} Current hour and minute from the DS1307 register.
    \item \textbf{Decision Logic:} The DS1307 continues tracking time using its backup battery during power loss. When power is restored to the ESP32-CAM, the board reads the current time from the RTC. If the time is within the scheduled night security window (e.g., 10:00 PM to 6:00 AM) and the system was previously armed, it automatically restores armed mode.
    \item \textbf{Local Alarm:} The system indicates recovery by flashing the red LED.
    \item \textbf{Cloud/Email/Telegram Notification:} Updates the cloud variable \texttt{system\_armed = true} and transmits a warning message: ``System recovered after power outage, security mode re-armed.''
    \item \textbf{Dashboard View:} The dashboard shows \texttt{system\_armed} as active and displays the recovery event in the logs.
    \item \textbf{Limitation/Assumption:} Assumes the DS1307 battery is charged and the RTC has been pre-configured with the correct time.
\end{itemize}
\subsection{Unknown WiFi Device Warning}
\begin{itemize}
    \item \textbf{Real-Life Scenario:} A suspect lingers outside the kitchen window, carrying a smartphone with active WiFi.
    \item \textbf{Sensors Involved:} ESP32-CAM WiFi module.
    \item \textbf{Data Read:} Nearby WiFi probe request packets and source MAC addresses.
    \item \textbf{Decision Logic:} The ESP32-CAM enters Promiscuous Mode periodically to capture ambient WiFi packets. If an unknown MAC address (not in the whitelist) is detected repeatedly (e.g., more than 5 times) within a specific time window, it is classified as a lingering threat.
    \item \textbf{Local Alarm:} None.
    \item \textbf{Cloud/Email/Telegram Notification:} Sends a warning notification to the Telegram chat group: ``Unknown WiFi device lingering near kitchen.''
    \item \textbf{Dashboard View:} The admin dashboard displays the detected MAC address and increments \texttt{unknown\_wifi\_count}.
    \item \textbf{Limitation/Assumption:} Deployed as an extended/partially implemented feature. Random MAC addresses generated by modern mobile operating systems can cause false positives.
\end{itemize}
\subsection{Automatic Arm/Disarm}
\begin{itemize}
    \item \textbf{Real-Life Scenario:} All family members leave the house for work during the day, forgetting to arm the system.
    \item \textbf{Sensors Involved:} ESP32-CAM WiFi module, RFID RC522 reader.
    \item \textbf{Data Read:} Whitelisted MAC addresses and RFID card UID readings.
    \item \textbf{Decision Logic:} The ESP32-CAM scans for the MAC addresses of whitelisted family phones. If all whitelisted devices are absent for more than 15 minutes, the system automatically arms itself. When a whitelisted device is detected again or a valid RFID card is swiped, the system disarms.
    \item \textbf{Local Alarm:} The buzzer emits a double beep to confirm disarming or a single long beep to indicate arming.
    \item \textbf{Cloud/Email/Telegram Notification:} Automatically updates \texttt{system\_armed} to reflect the current security state.
    \item \textbf{Dashboard View:} Users see the arming status toggle between active and inactive.
    \item \textbf{Limitation/Assumption:} Deployed as an extended/partially implemented feature. If a family member's phone battery dies, the system may arm while they are still home.
\end{itemize}
\subsection{Anti-Tampering Detection}
\begin{itemize}
    \item \textbf{Real-Life Scenario:} An intruder covers the LDR light sensor with black tape or a cloth to disable detection.
    \item \textbf{Sensors Involved:} LDR light sensor and PIR motion sensor.
    \item \textbf{Data Read:} Raw LDR light intensity (\texttt{ldr\_value}), LDR light level delta (\texttt{ldr\_delta}), and PIR motion status (\texttt{pir\_detected}).
    \item \textbf{Decision Logic:} If the LDR value drops suddenly (resulting in a high \texttt{ldr\_delta} value exceeding a set threshold) while the PIR detects motion in front of the device, the system identifies it as a tampering attempt.
    \item \textbf{Local Alarm:} Immediately activates the red LED and sounds the buzzer.
    \item \textbf{Cloud/Email/Telegram Notification:} Sets \texttt{sabotage\_alert = true}, updates the status to \texttt{SABOTAGE\_ALERT}, sends an email alert, and dispatches a Telegram message.
    \item \textbf{Dashboard View:} Displays a critical red alert indicator for sabotage on the admin and parent dashboards.
    \item \textbf{Limitation/Assumption:} Rapid changes in ambient light (e.g., cloud cover during the day or power outages turning off kitchen lights) can simulate a tampering event if not properly calibrated.
\end{itemize}
\subsection{SOS Emergency Alert from Arduino Cloud}
\begin{itemize}
    \item \textbf{Real-Life Scenario:} A child home alone hears a noise in the kitchen and triggers an emergency alarm from their phone.
    \item \textbf{Sensors Involved:} Arduino IoT Remote App, ESP32-CAM, Google Apps Script.
    \item \textbf{Data Read:} Cloud variables \texttt{sos\_child} and \texttt{sos\_adult}.
    \item \textbf{Decision Logic:} The ESP32-CAM monitors these cloud variables. If either variable changes to true, the board calculates the \texttt{sos\_level}:
    \begin{itemize}
        \item \emph{Level 1:} Local alarm only (Buzzer and LED).
        \item \emph{Level 2:} Local alarm and email notification to parent contacts.
        \item \emph{Level 3:} Local alarm and email notification with an escalation link to a simulated police contact email (no actual emergency services are contacted).
    \end{itemize}
    \item \textbf{Local Alarm:} Triggers continuous buzzer tones and flashes the red LED.
    \item \textbf{Cloud/Email/Telegram Notification:} Sets \texttt{alarm\_status = "SOS\_CHILD\_ALERT"} and calls Google Apps Script to dispatch the emails.
    \item \textbf{Dashboard View:} Shows active SOS status, displays the calculated \texttt{sos\_level}, and shows the generated \texttt{sos\_message}.
    \item \textbf{Limitation/Assumption:} Requires active internet connectivity to receive the command and send notifications.
\end{itemize}
% ==========================================
% SECTION 5: RESULTS AND DISCUSSION
% ==========================================
\chapter{Results and Discussion}
\section{Prototype Implementation}
The prototype system was assembled on a breadboard to verify the hardware and software design. The ESP32-CAM was mounted on an ESP32-CAM-MB programmer module to simplify programming and power delivery. The sensors were placed to maximize coverage: the PIR sensor was angled downward to monitor the room, the ultrasonic sensor was positioned parallel to the floor at a height of 80 cm to detect entry and filter pets, and the LDR was mounted on top of the enclosure to detect light changes. The completed assembly was housed in a compact plastic enclosure with cutouts for the sensors.
\section{Functional Test Results}
To validate the system, we tested each security function under simulated conditions. Table \ref{tab:test_results} summarizes the test results.
\begin{table}[H]
\centering
\caption{Prototype Functional Test Matrix}
\label{tab:test_results}
\begin{tabular}{lllp{6cm}}
    \toprule
    \textbf{Functionality} & \textbf{Status} & \textbf{Observed Output} & \textbf{Notes} \\
    \midrule
    Multi-sensor detection & PASS & Local alarms triggered; cloud state updated; email sent. & Required calibration of LDR threshold for daytime/nighttime transitions. \\
    Pet false-alarm filtering & DESIGN VALIDATED & Low-height objects (below 30 cm) did not trigger the alarm. & Verified using a simulated pet model (height: 25 cm). \\
    Power-loss recovery & PASS & Re-armed automatically during the scheduled security window. & The DS1307 RTC successfully maintained time accuracy. \\
    Unknown WiFi warning & PARTIAL & MAC addresses captured; alerts sent. & Random MAC address generation occasionally triggered false alerts. \\
    Automatic arm/disarm & PARTIAL & System armed/disarmed based on MAC detection. & Had a slight delay in detecting when devices left the network area. \\
    Anti-tampering detection & PASS & Triggered immediately when LDR was covered while PIR active. & Effectively detected attempts to block the sensor array. \\
    SOS emergency alert & PASS & Cloud button inputs triggered local alarms and dispatched emails. & Tested all 3 SOS levels successfully with simulated emails. \\
    \bottomrule
\end{tabular}
\end{table}
\section{Discussion}
The prototype testing confirmed that the multi-sensor approach provides a more reliable security solution than single-sensor designs. Using PIR and ultrasonic sensors together significantly reduced false alarms from small movements and low-profile objects, showing the value of sensor fusion. The anti-tampering function successfully identified attempts to block the LDR, adding a layer of physical security.
The DS1307 RTC proved reliable for power-loss recovery, ensuring the system restored its armed state once power returned. The integration of Google Apps Script and Telegram worked well, delivering remote alerts without overload. However, the WiFi-based features (unknown device detection and automatic arm/disarm) are still experimental due to WiFi scan delays and MAC randomization on modern smartphones.
% ==========================================
% SECTION 6: LIMITATIONS AND FUTURE WORK
% ==========================================
\chapter{Limitations and Future Work}
\section{System Limitations}
\begin{itemize}
    \item \textbf{Fire Detection Scope:} Fire detection was considered in earlier documentation but was removed from the final report scope to keep this version focused on anti-theft and emergency alert features. Therefore, the system cannot detect smoke or fires.
    \item \textbf{Network Dependency:} Remote notifications (emails, Telegram messages, and cloud updates) require an active internet connection. If the local router loses power or internet connectivity, the system can only sound local alarms.
    \item \textbf{WiFi Detection Reliability:} The unknown WiFi device warning and automatic arm/disarm features are prone to inconsistencies due to MAC randomization on modern smartphones and network scan delays.
    \item \textbf{Power Consumption:} The ESP32-CAM requires a constant power connection to run its WiFi transceiver. It cannot operate long on small backup batteries without deep-sleep optimizations.
\end{itemize}
\section{Future Work}
\begin{itemize}
    \item \textbf{Transition to a Custom PCB:} Designing a custom printed circuit board (PCB) would improve system reliability by replacing loose jumper wires with solid, soldered connections.
    \item \textbf{Edge Image Processing:} Future versions could use the ESP32-CAM's camera to run lightweight on-board face detection or motion capture, reducing false positives further and capturing images of intruders.
    \item \textbf{Battery Backup System:} Adding a rechargeable lithium-ion battery backup circuit would allow the system to keep monitoring locally and sending alerts during power outages.
    \item \textbf{Alternative Wireless Protocols:} Replacing or supplementing WiFi with low-power, long-range protocols like Zigbee or LoRa would improve battery life and reliability.
\end{itemize}
% ==========================================
% SECTION 7: CONCLUSION
% ==========================================
\chapter{Conclusion}
This project successfully demonstrated the design and implementation of an IoT-Based Multi-Sensor Anti-Theft System for Kitchen Areas. By combining low-cost sensors (PIR, ultrasonic, LDR) and using sensor fusion logic on the ESP32-CAM, the system provides reliable security monitoring while minimizing false alarms from pets. The DS1307 RTC ensures the system recovers its security state automatically after power outages. 
The integration with Arduino IoT Cloud, Google Apps Script, and Telegram Bot delivers a robust notification network with role-based access for different family members. While there are still limitations in the experimental WiFi features, the prototype shows the viability of multi-sensor fusion for accessible, reliable smart home security.
% ==========================================
% REFERENCES
% ==========================================
\bibliographystyle{plain}
\bibliography{references}
\end{document}
\begin{figure}
    \centering
    \includegraphics[width=0.5\linewidth]{images.png}
    \caption{Enter Caption}
    \label{fig:placeholder}
\end{figure}